%==============================================================================
%  SynectivSec: A Composition Theorem for Federated Agent Authority
%
%  Author: Abhijith S (ORCID 0009-0009-9232-7820)
%  Artifact: https://github.com/kakashi-kx/synectivsec
%
%  Compile with: pdflatex synectivsec.tex (run twice for refs)
%==============================================================================
\documentclass[11pt,a4paper]{article}

% --- Encoding & fonts -------------------------------------------------------
\usepackage[utf8]{inputenc}
\usepackage[T1]{fontenc}
\usepackage{lmodern}
\usepackage{microtype}

% --- Layout -----------------------------------------------------------------
\usepackage[margin=1in]{geometry}
\usepackage{parskip}

% --- Math -------------------------------------------------------------------
\usepackage{amsmath}
\usepackage{amssymb}
\usepackage{amsthm}
\usepackage{stmaryrd}      % for \llbracket, \rrbracket
\usepackage{mathrsfs}

% --- Theorem environments ---------------------------------------------------
\theoremstyle{plain}
\newtheorem{theorem}{Theorem}
\newtheorem{lemma}[theorem]{Lemma}
\newtheorem{corollary}[theorem]{Corollary}
\newtheorem{proposition}[theorem]{Proposition}

\theoremstyle{definition}
\newtheorem{definition}[theorem]{Definition}
\newtheorem{example}[theorem]{Example}

\theoremstyle{remark}
\newtheorem{remark}[theorem]{Remark}

% --- Tables -----------------------------------------------------------------
\usepackage{booktabs}

% --- Hyperlinks -------------------------------------------------------------
\usepackage[hidelinks,colorlinks=true,linkcolor=blue!50!black,
            citecolor=blue!50!black,urlcolor=blue!50!black]{hyperref}
\usepackage{cleveref}

% --- PDF metadata -----------------------------------------------------------
\hypersetup{
  pdftitle={SynectivSec: A Composition Theorem for Federated Agent Authority},
  pdfauthor={Abhijith S},
  pdfsubject={Composition theorem for delegation chains, sequence policies, and federation},
  pdfkeywords={agent security, delegation chains, sequence policies, federation,
               TLA+, composition theorem, independence result},
}

% --- Code listings ----------------------------------------------------------
\usepackage{listings}
\usepackage{xcolor}
\lstdefinelanguage{TLA+}{
  morekeywords={MODULE,EXTENDS,CONSTANTS,VARIABLES,ASSUME,THEOREM,LEMMA,
                INSTANCE,LOCAL,CONSTANT,VARIABLE,DEFINE,RECURSIVE,
                CASE,IF,THEN,ELSE,LET,IN,CHOOSE,DOMAIN,SUBSET,UNION,
                EXCEPT,ENABLED,UNCHANGED,WF_,SF_,\A,\E,\in,\notin,
                \subseteq,\subset,\cup,\cap,\o,\X,\div,\mod,\leq,\geq,
                \neq,=,<,>,=>,<=>,~,/\,\/,[] ,<>,\EE,\AA},
  sensitive=true,
  morecomment=[l]{\*},
  morestring=[b]"
}
\lstset{
  basicstyle=\ttfamily\footnotesize,
  keywordstyle=\color{blue!60!black},
  commentstyle=\color{green!50!black}\itshape,
  stringstyle=\color{red!60!black},
  frame=single,
  breaklines=true,
  showstringspaces=false,
  columns=fullflexible
}

% --- Bibliography -----------------------------------------------------------
% We use a simple manual bibliography to avoid requiring .bib tooling.
% References are listed in \begin{thebibliography} at the end.

% --- Title ------------------------------------------------------------------
\title{SynectivSec: A Composition Theorem for\\Federated Agent Authority}

\author{
  Abhijith S\thanks{%
    Independent Researcher.\\
    ORCID: \href{https://orcid.org/0009-0009-9232-7820}{0009-0009-9232-7820}.\\
    GitHub: \href{https://github.com/kakashi-kx}{@kakashi-kx}.\\
    Artifact: \href{https://github.com/kakashi-kx/synectivsec}{github.com/kakashi-kx/synectivsec}.
  }
}

\date{September 2026}

%==============================================================================
\begin{document}
\maketitle

%------------------------------------------------------------------ Abstract --
\begin{abstract}
Delegation chains, sequence policies, and cross-domain federation each have
formal treatments in the security and distributed-systems literature, but no
result states when their \emph{composition} preserves the safety properties of
its parts. We give three compatibility conditions under which it does, prove
the preservation theorem by structural induction over the transition relation,
and show that dropping any one condition admits a concrete violation. We
model-check the construction in TLA+ on a bounded instance
(11{,}185{,}890 states generated, 300{,}447 distinct, zero violations of the
five stated invariants) and state exactly where the proof and the verification
stop. The individual conditions are not novel---each restates prior work---but
the composition operator, the joint formulation of the conditions, and the
preservation theorem are. The artifact is released at
\url{https://github.com/kakashi-kx/synectivsec}.
\end{abstract}

%---------------------------------------------------------------- Introduction --
\section{Introduction}
\label{sec:intro}

Autonomous AI agents take actions: they call tools, move funds, send
communications, and invoke other agents. Each action is the composition of
several layers of authority and policy. A delegation chain establishes how
authority attenuates from a root principal to the acting agent. A sequence
policy constrains the order and multiplicity of actions over time. A
federation graph determines which trust domains accept whose attestations and
how revocation propagates across organizational boundaries.

Each of these layers has been studied separately. Delegation chains are
formalized in the access-control literature as far back as Lampson, Abadi,
Burrows, and Wobber~\cite{lampson1992authentication} and in the SPKI/SDSI
certificate theory~\cite{rfc2693spki}; more recently, in the delegation-chain
calculi of the agent-security literature~\cite{sentinelagent2026}. Sequence
policies over action histories are grounded in
Schneider's work on enforceable security policies~\cite{schneider2000enforceable}
and in modern agent frameworks such as the Agentic Principal
Chain~\cite{boundedagents2026}. Cross-domain federation is discussed in
federated governance testbeds~\cite{wang2026federated} and in the meta-framework
COA-MAS v2~\cite{coamas2026}.

What is missing is a \emph{composition theorem}: a statement of the conditions
under which the composition of all three layers preserves the safety invariants
of each. This absence is not accidental. Compositionality of security
properties is difficult in general: McLean~\cite{mclean1994general} showed that
safety is not preserved under composition without strong assumptions, and
recent work in agent security has made the difficulty concrete.
Spera~\cite{spera2026noncompositionality} proves a non-compositionality result
showing that two individually safe agents can jointly reach a forbidden state
through an emergent conjunctive dependency. AgentRFC~\cite{agentrfc2026}
documents 20 cross-protocol violations across five composition patterns.

The difficulty is real, but the composition question is not hopeless. We give
a positive result for a specific class of systems: those satisfying three
compatibility conditions. Under these conditions, the composition preserves
five safety invariants. We prove the theorem by structural induction over the
transition relation and verify it in TLA+ on a bounded instance. We further
show that each condition is non-redundant: dropping any one admits a
counterexample.

\paragraph{Contributions.}

\begin{enumerate}
  \item \textbf{The composition operator $\otimes$} for the joint treatment of
  delegation chains, sequence policies, and federated attestation
  (\S\ref{sec:model}).
  \item \textbf{Three compatibility conditions}---Policy Compatibility (PC),
  Revocation Synchronization (RS), and Cross-Boundary Attestation
  Transitivity (CBAT)---as a sufficient set for the preservation of five
  safety invariants (\S\ref{sec:conditions}).
  \item \textbf{A preservation theorem} (\Cref{thm:preservation}) proved by
  structural induction over the transition relation
  (\S\ref{sec:proof}).
  \item \textbf{A mechanized verification} in TLA+ on a bounded instance, with
  the raw output and reproduction instructions (\S\ref{sec:verification}).
  \item \textbf{An independence result}: each condition is non-redundant, with
  a runnable counterexample per condition (\S\ref{sec:independence}).
\end{enumerate}

\paragraph{What is not claimed.}

We do not claim novelty for any of the three conditions in isolation. CBAT
restates the \emph{speaks-for} transitivity of
Lampson~et~al.~\cite{lampson1992authentication} with the bounded variant the
same paper suggests. RS applies the validity-interval concept of SPKI
RFC~2693~\cite{rfc2693spki} to a domain graph with an explicit propagation
bound. PC relates to the sequence-policy work of
Bounded Agents~\cite{boundedagents2026} and
Schneider~\cite{schneider2000enforceable}. The claimed contribution is the
\emph{composition} of the three conditions, not their individual content. A
full lineage is given in \Cref{tab:lineage}.

\paragraph{Paper organization.}

\Cref{sec:adversaries} presents the adversary model that motivates the
composition question. \Cref{sec:model} defines the three layers and the
composition operator. \Cref{sec:conditions} states the three compatibility
conditions and the preservation theorem, with a per-condition lineage table.
\Cref{sec:proof} gives the proof. \Cref{sec:verification} reports the TLA+
verification. \Cref{sec:independence} reports the independence result and the
attack harness. \Cref{sec:related} discusses related work.
\Cref{sec:limitations} states limitations and threats to validity.
\Cref{sec:conclusion} concludes. The appendices contain the full proof
(\Cref{app:proof}) and the TLA+ module (\Cref{app:tla}).

%----------------------------------------------------- Composition adversaries --
\section{Composition adversaries}
\label{sec:adversaries}

The adversaries that motivate this work are not the classical single-layer
attackers. They exist only because the three layers compose. We catalogue them
by the composition-specific capability each requires, following the attack
catalog in the artifact's \texttt{docs/ALGEBRA.md} and
\texttt{redteam/README.md}.

\begin{definition}[Composition adversary]
A \emph{composition adversary} is an attacker whose power derives from the
interaction of the delegation, sequence, and federation layers, and whose goal
is to cause an action to occur that no single layer would authorize.
\end{definition}

\Cref{tab:adversaries} lists the eight adversary classes we consider. Each is
characterized by the composition-specific capability it requires and the
invariant it seeks to violate.

\begin{table}[t]
\centering
\small
\begin{tabular}{@{}p{0.8cm}p{3.6cm}p{4.2cm}p{4.2cm}@{}}
\toprule
\textbf{ID} & \textbf{Capability} & \textbf{Goal} & \textbf{Invariant at risk} \\
\midrule
AC-1 & Holds a valid delegation chain & Extend it beyond intent, relying on the sequence policy not seeing the extension & CI2, CI4 \\
AC-2 & Controls two chains & Interleave actions such that no single chain's sequence policy fires & CI2 \\
AC-3 & Holds attestations in two domains & Launder authority across non-transitive attestation paths & CI4 \\
AC-4 & Controls timing across domains & Act in the window before revocation propagates & CI3 \\
AC-5 & Can craft action parameters & Produce an action outside the sequence policy's recognized alphabet & CI2 \\
AC-6 & Observes one domain's attestations & Replay an attestation in a domain where it was not issued & CI4 \\
AC-7 & Controls multiple agents & Jointly reach a forbidden state, each individually authorized & CI1--CI5 \\
AC-8 & Exploits trust-boundary ambiguity & Cause two domains to each believe the other's policy applies & CI5 \\
\bottomrule
\end{tabular}
\caption{Composition adversary classes. Each exists only under composition.}
\label{tab:adversaries}
\end{table}

The key observation is that AC-1 through AC-8 are not hypothetical. Each
corresponds to an attack that has been demonstrated in the agent-security
literature or is a natural extension of one. AC-7 (emergent conjunctive
authority) is the Spera non-compositionality
result~\cite{spera2026noncompositionality}. AC-3 and AC-6 (attestation
laundering and replay) are natural cross-domain analogues of attacks documented
in the trust-management literature. AC-4 (revocation race) is the composition
of the classic time-of-check/time-of-use problem with cross-domain
propagation. AC-5 (alphabet escape) is specific to the composition of the
sequence layer with delegation.

Against this adversary set, the composition question is: under what conditions
does the composition of the three layers prevent all of AC-1 through AC-8?

%------------------------------------------------------------------ Model --
\section{Model}
\label{sec:model}

We define the three layers, the composition operator, and the safety
invariants. The definitions here are the specification; the artifact's
\texttt{docs/ALGEBRA.md} mirrors them, and
\texttt{tla/Composition.tla} implements them in TLA+.

\subsection{Principals and capabilities}

A \emph{principal} is a named entity that can sign statements, be delegated
authority, and perform actions. Write $\mathsf{Principal}$ for the finite set
of principals in a system, and $p, q, r, \ldots$ for elements of it.

A \emph{capability} is an atomic permission. Write $\mathsf{Capability}$ for
the finite set of capabilities, and $c, c', \ldots$ for elements. Capabilities
form a lattice $\langle \mathsf{Capability}, \sqsubseteq \rangle$ in which
$c \sqsubseteq c'$ means $c$ is at least as strong as $c'$ (i.e., $c$ permits
everything $c'$ permits, and possibly more). Each principal $p$ has a
capability set $\kappa(p) \subseteq \mathsf{Capability}$.

\subsection{Delegation chains}

A \emph{delegation chain} is a finite sequence
\begin{equation*}
D = \langle \kappa_0, \sigma_0, \kappa_1, \sigma_1, \ldots, \kappa_n, \sigma_n \rangle
\end{equation*}
where each $\kappa_i$ is a capability set, each $\sigma_i$ is a signature by
the holder of $\kappa_{i-1}$ attesting to $\kappa_i$, and the attenuation
condition holds:
\begin{equation*}
\kappa_n \sqsubseteq \kappa_{n-1} \sqsubseteq \cdots \sqsubseteq \kappa_0.
\end{equation*}
The root capability set $\kappa_0$ is granted by a designated root principal.
Let $\mathsf{Reach}(D)$ be the set of actions producible under $D$ without
violating attenuation.

\paragraph{Prior art.} The delegation-chain model is standard. Lampson, Abadi,
Burrows, and Wobber~\cite{lampson1992authentication} introduce the
\emph{speaks-for} relation and the handoff axiom in \S3.2 and \S3.3
respectively. SPKI RFC~2693~\cite{rfc2693spki} formalizes certificate-chain
reduction in \S6.3. SentinelAgent~\cite{sentinelagent2026} provides a modern
delegation chain calculus with TLA+ verification. The definition above restates
the standard model; it does not introduce new primitives.

\subsection{Sequence policies}

A \emph{sequence context} is a pair $S = \langle \Sigma, \varphi \rangle$
where $\Sigma$ is a finite \emph{action alphabet} and $\varphi : \Sigma^* \to
\{\top, \bot\}$ is a decidable predicate on finite action histories. We write
$h \models \varphi$ for $\varphi(h) = \top$.

We assume $\varphi$ is \emph{prefix-closed}: $h \models \varphi$ implies
$\mathrm{prefix}(h, i) \models \varphi$ for all $i \le |h|$. Prefix-closure
characterizes the safety policies in the sense of
Alpern and Schneider~\cite{alpern1985defining}. Liveness policies are outside
the scope of this paper.

Let $\mathsf{Reach}(S)$ be the set of histories $h \in \Sigma^*$ with $h
\models \varphi$.

\paragraph{Prior art.} Runtime enforcement of safety policies is
Schneider's~\cite{schneider2000enforceable}. The specific formalism of
sequence policies over agent actions is from Bounded
Agents~\cite{boundedagents2026} (the Agentic Principal Chain). The definition
above restates the standard model.

\subsection{Trust domains and federation}

A \emph{trust domain} is a named region of policy and trust authority, written
$T$ or $T_i$. Let $\mathsf{Domain}$ be the set of all domains. A
\emph{federation graph} is a tuple
\begin{equation*}
F = \langle \mathcal{T}, \mathcal{A}, \delta \rangle
\end{equation*}
where:
\begin{itemize}
  \item $\mathcal{T} \subseteq \mathsf{Domain}$ is a finite set of domains;
  \item $\mathcal{A} \subseteq \mathcal{T} \times \mathcal{T}$ is the
  \emph{attestation relation}, with $\mathcal{A}(T_i, T_j)$ meaning ``$T_i$
  attests the claims of $T_j$'';
  \item $\delta \in \mathbb{N}$ is the \emph{revocation propagation bound}.
\end{itemize}

Let $\mathsf{Reach}(F, T)$ be the set of domains reachable from $T$ via
$\mathcal{A}$, including $T$ itself. We define $\mathsf{Reach}$ as the
transitive closure of $\mathcal{A}$ together with the reflexive base case:
\begin{equation*}
\mathsf{Reach}(F, T) \;=\; \{T\} \;\cup\; \{T' : \mathcal{A}(T, T')\}
\;\cup\; \{T'' : \exists T'. \mathcal{A}(T, T') \wedge \mathcal{A}(T', T'')\}
\;\cup\; \cdots
\end{equation*}

The revocation propagation bound $\delta$ is discussed in
\S\ref{sec:conditions}.

\paragraph{Prior art.} Cross-domain trust is discussed in Wang's federated
governance testbed~\cite{wang2026federated} and COA-MAS
v2~\cite{coamas2026}. Attestation transitivity is from Lampson et
al.~\cite{lampson1992authentication} \S3.2 and SPKI RFC~2693~\cite{rfc2693spki}
\S6.3. The definition above restates the standard model.

\subsection{Composition states}

A \emph{composition state} is a tuple
\begin{equation*}
s = \langle d, h, T, r, t \rangle
\end{equation*}
where:
\begin{itemize}
  \item $d \in \mathsf{Reach}(D)$ is the current position in a delegation chain;
  \item $h \in \Sigma^*$ is the current action history;
  \item $T \in \mathcal{T}$ is the current trust domain;
  \item $r : \mathsf{Principal} \to \mathbb{N} \cup \{\infty\}$ maps each
  principal to its revocation time (or $\infty$ if never revoked);
  \item $t \in \mathbb{N}$ is the current logical time.
\end{itemize}

Write $\mathsf{States}(D, S, F)$ for the set of all composition states.

\subsection{The transition relation}

The transition relation $\to$ on composition states is given by three rules.

\paragraph{Action transition.} A principal $p$ performs an action $a \in
\Sigma$ in domain $T$:
\begin{equation*}
\langle d, h, T, r, t \rangle \;\to\; \langle d, h \frown \langle \langle p, a, t \rangle \rangle, T, r, t + 1 \rangle
\end{equation*}
subject to five guards:
\begin{align*}
\text{G1 (Reachability):} & \quad T \in \mathsf{Reach}(F, \mathsf{RootDomain}(p)) \\
\text{G2 (Authorization):} & \quad a \in \kappa(p) \;\wedge\; t < r[p] \\
\text{G3 (Sequence):} & \quad h \frown \langle a \rangle \models \varphi \\
\text{G4 (Length):} & \quad |h| < \mathit{MaxHistory} \\
\text{G5 (Boundary):} & \quad \mathsf{DomainPolicy}(T, a) = \top
\end{align*}

\paragraph{Revocation transition.} A principal $p$ is revoked at time $t$:
\begin{equation*}
\langle d, h, T, r, t \rangle \;\to\; \langle d, h, T, r[p \mapsto t], t + 1 \rangle
\end{equation*}
subject to $r[p] = \infty$ (not already revoked).

\paragraph{Domain transition.} A principal $p$ moves to a domain $T' \in
\mathcal{T}$:
\begin{equation*}
\langle d, h, T, r, t \rangle \;\to\; \langle d, h, T', r, t + 1 \rangle
\end{equation*}
subject to $T' \in \mathsf{Reach}(F, \mathsf{RootDomain}(p))$.

The TLA+ model in \texttt{tla/Composition.tla} implements these rules as
\texttt{DoAction}, \texttt{Revoke}, and \texttt{CrossDomain} respectively.

\subsection{The composition operator $\otimes$}

\begin{definition}[Composition]
The composition $D \otimes S \otimes F$ is the labelled transition system
$\langle \mathsf{States}, \to, s_0 \rangle$ where:
\begin{itemize}
  \item $\mathsf{States} = \mathsf{Reach}(D) \times \Sigma^* \times \mathcal{T}$
  (with the revocation map and time implicit in the execution);
  \item $\to$ is the transition relation from \S\ref{sec:model};
  \item $s_0 = \langle d_0, \varepsilon, T_0, r_0, 0 \rangle$ is the initial
  state, with $r_0[p] = \infty$ for all $p$.
\end{itemize}
The \emph{reachable set} $\mathsf{Reach}(D \otimes S \otimes F)$ is the
transitive closure of $\to$ from $s_0$.
\end{definition}

The operator $\otimes$ is binary and associative by construction:
$(D \otimes S) \otimes F = D \otimes (S \otimes F)$. It is \emph{not
commutative} in general: $D \otimes S \neq S \otimes D$, because the sequence
policy constrains actions, and actions are only produced under delegation.
This order-sensitivity mirrors the Composition Safety principle of
AgentRFC~\cite{agentrfc2026}.

\subsection{Safety invariants}

We state five invariants on composition states. Write $h[i] = \langle p_i,
a_i, t_i \rangle$ for the $i$-th entry of the history.

\begin{definition}[CI1--CI5]
For a state $s = \langle d, h, T, r, t \rangle$:
\begin{align*}
\mathsf{CI1}(s) &\;=\; \forall i \in 1..|h|.\; \exists p \in \mathsf{Principal}.\; h[i].\mathit{action} \in \kappa(p) \\
\mathsf{CI2}(s) &\;=\; h \models \varphi \\
\mathsf{CI3}(s) &\;=\; \forall i \in 1..|h|.\; h[i].\mathit{time} < r[h[i].\mathit{principal}] \\
\mathsf{CI4}(s) &\;=\; T \in \mathsf{Reach}(F, \mathsf{RootDomain}(\mathit{currentPrincipal})) \\
\mathsf{CI5}(s) &\;=\; \forall T \in \mathcal{T}, a \in \Sigma.\;
                  \mathsf{DomainPolicy}(T, a) = \top \vee \mathsf{DomainPolicy}(T, a) = \bot
\end{align*}
\end{definition}

Informally:
\begin{itemize}
  \item \textbf{CI1 (Authority containment)}: every action in history is
  authorized by some principal's capability set.
  \item \textbf{CI2 (Sequence soundness)}: the action history satisfies the
  sequence policy.
  \item \textbf{CI3 (Revocation freshness)}: no action follows its actor's
  revocation.
  \item \textbf{CI4 (Attestation soundness)}: the current principal acts only
  in domains reachable from its root domain.
  \item \textbf{CI5 (Boundary determinism)}: the applicable domain policy is
  defined at every state.
\end{itemize}

%------------------------------------------------- Conditions and theorem --
\section{Compatibility conditions and the preservation theorem}
\label{sec:conditions}

Three compatibility conditions, taken jointly, ensure the preservation of the
five invariants. Each condition is stated below, followed by its lineage to
prior work.

\subsection{The conditions}

\begin{definition}[Policy Compatibility, PC]
\label{def:pc}
The composition $D \otimes S \otimes F$ satisfies \emph{Policy Compatibility}
if every action producible under any chain in $D$ is classified by the
sequence policy's alphabet:
\begin{equation*}
\mathsf{PC}(D, S) \;=\; \forall d \in \mathsf{Reach}(D).\; \mathsf{Actions}(d) \subseteq \Sigma.
\end{equation*}
\end{definition}

PC ensures that the sequence policy sees every action that can occur. Without
it, an action outside $\Sigma$ escapes the sequence constraint entirely.

\begin{definition}[Revocation Synchronization, RS]
\label{def:rs}
The composition satisfies \emph{Revocation Synchronization} with bound
$\delta$ if revocation in one domain is effective in every domain reachable by
an attestation path, within $\delta$ time units:
\begin{equation*}
\mathsf{RS}(F, \delta) \;=\; \forall T_i, T_j \in \mathcal{T}.\;
\mathsf{path}(\mathcal{A}, T_i, T_j) \;\Rightarrow\;
\mathsf{revoke}(T_i, t) \;\Rightarrow\; \mathsf{effective}(T_j, t + \delta).
\end{equation*}
\end{definition}

RS ensures that a revoked principal cannot continue to act in a domain that
has not yet seen the revocation. Without it, the propagation window admits a
race.

\begin{definition}[Cross-Boundary Attestation Transitivity, CBAT]
\label{def:cbat}
The composition satisfies \emph{Cross-Boundary Attestation Transitivity} if
attestation is transitive across domains unless explicitly bounded:
\begin{equation*}
\mathsf{CBAT}(F) \;=\; \forall T_i, T_j, T_k \in \mathcal{T}.\;
\mathcal{A}(T_i, T_j) \wedge \mathcal{A}(T_j, T_k) \;\Rightarrow\;
\mathcal{A}(T_i, T_k) \vee \mathsf{explicitly\_bounded}(T_i, T_k).
\end{equation*}
\end{definition}

CBAT ensures that reachability is closed under composition of attestations.
Without it, an agent can act in a domain reachable only through an
intermediate domain whose attestation was not transitively closed---an
attestation laundering attack.

\subsection{Lineage of the conditions}
\label{sec:lineage}

None of the three conditions is novel in isolation. \Cref{tab:lineage} states,
for each condition, the prior work it restates and the aspect (if any) that
extends or refines that prior work.

\begin{table}[t]
\centering
\small
\begin{tabular}{@{}p{0.7cm}p{6.4cm}p{5.8cm}@{}}
\toprule
\textbf{Cond.} & \textbf{Restates} & \textbf{Where} \\
\midrule
PC & Scope/budget tracking across action sequences; runtime enforcement of safety policies & Bounded Agents~\cite{boundedagents2026}; Schneider~\cite{schneider2000enforceable} \\
RS & Validity intervals on certificates; the question ``how long are you willing to let the world believe and act on a statement you know to be false?'' & SPKI RFC~2693~\cite{rfc2693spki}, \S5.4 \\
CBAT & Speaks-for transitivity: ``it is also easy to show that $\vdash$ is monotonic in both arguments and that $\Rightarrow$ is transitive'' & Lampson et al.~\cite{lampson1992authentication}, \S3.2 \\
CBAT (bounding) & ``the general axiom is too powerful\ldots if the conclusion uses a qualified form of $\Rightarrow$ it may be more acceptable''; closest formal analog is RT's linked roles & Lampson et al.~\cite{lampson1992authentication}, \S3.3 (note on P10); Li, Mitchell, Winsborough~\cite{rt2002} \\
\bottomrule
\end{tabular}
\caption{Prior-art lineage of the three compatibility conditions. Each condition
restates prior work; the contribution is the joint treatment.}
\label{tab:lineage}
\end{table}

\subsection{The preservation theorem}

\begin{theorem}[Composition Preservation]
\label{thm:preservation}
If $\mathsf{PC}(D, S) \wedge \mathsf{RS}(F, \delta) \wedge \mathsf{CBAT}(F)$
hold, then for every state $s \in \mathsf{Reach}(D \otimes S \otimes F)$ and
every $k \in \{1, \ldots, 5\}$:
\begin{equation*}
\mathsf{CI}k(s).
\end{equation*}
\end{theorem}

The contrapositive gives a concrete attack surface per condition:

\begin{theorem}[Emergent Authority, contrapositive]
\label{thm:emergent}
If $\neg\mathsf{PC} \vee \neg\mathsf{RS} \vee \neg\mathsf{CBAT}$, there exists
a trajectory $\tau$ in $\mathsf{Reach}(D \otimes S \otimes F)$ that violates
at least one of $\mathsf{CI1}, \ldots, \mathsf{CI5}$.
\end{theorem}

\Cref{thm:emergent} is the red-teamer's reading of the theorem: every failure
of a compatibility condition yields a concrete attack. \Cref{sec:independence}
demonstrates this for each condition by a runnable counterexample.

%------------------------------------------------------------------- Proof --
\section{Proof of the preservation theorem}
\label{sec:proof}

The proof is by structural induction over the transition relation. The base
case establishes all five invariants at the initial state. The inductive step
shows that each of the three transition rules preserves all five invariants.
The full proof is given in \Cref{app:proof}; here we give the induction
outline and state where each compatibility condition is used.

\subsection{Proof outline}

\begin{proof}[Proof of \Cref{thm:preservation}, by structural induction]
\emph{Base case.} Let $s_0 = \langle d_0, \varepsilon, T_0, r_0, 0 \rangle$.
Since $h = \varepsilon$, the quantifiers in CI1, CI2, and CI3 are vacuous.
CI4 holds because $T_0 \in \mathsf{Reach}(F, T_0) =
\mathsf{Reach}(F, \mathsf{RootDomain}(p_0))$ by the reflexive base case of
$\mathsf{Reach}$. CI5 holds because $\mathsf{DomainPolicy}$ is total by
definition.

\emph{Inductive step.} Assume $\mathsf{CI1}(s) \wedge \cdots \wedge
\mathsf{CI5}(s)$ for some $s \in \mathsf{Reach}(D \otimes S \otimes F)$. Let
$s \to s'$. We consider the three transition rules in turn.

\emph{T-Action:} $s' = \langle d, h \frown \langle \langle p, a, t \rangle
\rangle, T, r, t + 1 \rangle$.
\begin{itemize}
\item CI1 follows from $\mathsf{CI1}(s)$ for old history entries and from
guard G2 for the new entry.
\item CI2 follows directly from guard G3.
\item CI3 follows from $\mathsf{CI3}(s)$ for old entries and from guard G2
($t < r[p]$) for the new entry. \textbf{RS is used here}: the value $r[p]$
must be consistent across domains, which is exactly what RS guarantees.
\item CI4 follows from guard G1. \textbf{CBAT is used here}: the reachability
relation $\mathsf{Reach}$ is defined as the transitive closure of
$\mathcal{A}$, and CBAT is what makes the transitive closure well-defined
across domains.
\item CI5 follows from $\mathsf{CI5}(s)$ and from guard G5. \textbf{PC is
used here}: PC ensures the new action is in $\Sigma$, so the domain policy
classification applies.
\end{itemize}

\emph{T-Revoke:} $s' = \langle d, h, T, r[p \mapsto t], t + 1 \rangle$. The
history, current principal, and current domain are unchanged. Only
$r$ changes, and only at $p$. The invariants CI1, CI2, CI4, CI5 are
immediate. For CI3: for $i$ with $h[i].\mathit{principal} \neq p$, the
inequality is preserved; for $i$ with $h[i].\mathit{principal} = p$, we have
$h[i].\mathit{time} < t$ because no action in $h$ is recorded at the current
time (actions increment $t$), so the inequality holds.

\emph{T-Cross:} $s' = \langle p, h, T', r, t + 1 \rangle$. The history is
unchanged. CI1, CI2, CI3 are preserved. CI4 follows from the transition
guard $T' \in \mathsf{Reach}(F, \mathsf{RootDomain}(p))$; \textbf{CBAT is
used here} for the same reason as in T-Action. CI5 is unchanged.
\end{proof}

\subsection{Where each condition is used}

\Cref{tab:conditions-use} summarizes the use of each compatibility condition
in the proof.

\begin{table}[t]
\centering
\small
\begin{tabular}{@{}p{1.2cm}p{5.5cm}p{7cm}@{}}
\toprule
\textbf{Cond.} & \textbf{Used in} & \textbf{Why} \\
\midrule
PC & CI2, CI5 preservation under T-Action & Ensures the new action is in $\Sigma$, so the sequence policy and domain policy classify it \\
RS & CI3 preservation under T-Action & Ensures revocation time $r[p]$ is consistent across domains, so the guard $t < r[p]$ is sound \\
CBAT & CI4 preservation under T-Action and T-Cross & Ensures $\mathsf{Reach}(F, T)$ is the transitive closure of $\mathcal{A}$, so the guard $T \in \mathsf{Reach}$ is well-defined \\
\bottomrule
\end{tabular}
\caption{Dependence of the proof on each compatibility condition.}
\label{tab:conditions-use}
\end{table}

The use is disjoint: each condition is used in a distinct part of the
preservation argument. This is what makes the conditions non-redundant, as
\Cref{sec:independence} demonstrates by counterexample.

%------------------------------------------------------ Mechanized verification --
\section{Mechanized verification}
\label{sec:verification}

The theorem is verified in TLA+ for a bounded instance. The verification is
not a proof of the general theorem---the proof in \S\ref{sec:proof} carries
that---but it catches specification bugs and confirms that the model is
coherent on a concrete configuration.

\subsection{The verified instance}

\begin{table}[t]
\centering
\small
\begin{tabular}{@{}ll@{}}
\toprule
\textbf{Parameter} & \textbf{Value} \\
\midrule
Principals & 6 ($R_1, PA, PB, R_2, PC, PD$) \\
Delegation chains & 2 (chain $D_1$ rooted in $T_1$; chain $D_2$ rooted in $T_2$) \\
Trust domains & 3 ($T_1, T_2, T_3$) \\
Capabilities & $\{\mathit{read}, \mathit{write}, \mathit{send}\}$ \\
Sequence policy & No two consecutive \texttt{send} actions \\
Attestation relation & $T_1 \to T_2$, $T_2 \to T_3$ \\
State constraint & $\mathit{time} \le 4$, $\mathit{Len}(\mathit{history}) \le 3$ \\
Invariants checked & CI1, CI2, CI3, CI4, CI5 \\
\bottomrule
\end{tabular}
\caption{The verified instance.}
\label{tab:instance}
\end{table}

The instance is deliberately minimal: it is small enough that TLC exhausts the
reachable state space in seconds, and rich enough that each invariant is
non-trivial. The sequence policy exercises two-chain interleaving. The
attestation relation exercises transitivity from $T_1$ to $T_3$ via $T_2$. The
action-level revocation tracking exercises CI3.

\subsection{Verification result}

TLC 2.19 exhausts the reachable state space with the following result:

\begin{verbatim}
Model checking completed. No error has been found.
11185890 states generated, 300447 distinct states found,
    0 states left on queue.
The depth of the complete state graph search is 5.
Finished in 33s.
\end{verbatim}

Zero violations of CI1--CI5 across all 300{,}447 distinct states. The
fingerprint collision probability reported by TLC is $1.0 \times 10^{-9}$
(based on actual fingerprints), so the exhaustive search is sound.

\subsection{What the verification does and does not show}

The verification establishes:
\begin{itemize}
  \item The five invariants are consistent with the transition relation on the
  stated instance.
  \item TLC's exhaustive search terminates with zero violations, so the
  theorem holds on this instance with high confidence.
\end{itemize}

The verification does \emph{not} establish:
\begin{itemize}
  \item The theorem for unbounded executions. The state constraint bounds
  $\mathit{time} \le 4$ and $\mathit{Len}(\mathit{history}) \le 3$.
  \item The theorem for larger instances. Scaling has been attempted but the
  state space grows beyond what TLC can handle; see
  \S\ref{sec:limitations}.
  \item The theorem in a proof assistant. The proof is by hand.
\end{itemize}

The verification complements the proof: the proof is general but by hand, the
verification is bounded but mechanical. Neither alone is sufficient.

%----------------------------------------------------------- Independence --
\section{Independence of the conditions}
\label{sec:independence}

The three conditions are individually non-redundant. Dropping any one while
keeping the other two admits a reachable state that violates a named
invariant. This is the independence of the hypothesis: no condition can be
removed without loss.

\paragraph{Method.} For each condition, the artifact's attack harness
constructs a composition state that (a) satisfies the other two conditions and
(b) violates the invariant tied to the dropped condition. The attack is
implemented as a guard-level weakening of the TLA+ model, followed by TLC.
The TLC trace is the witness.

\subsection{PC non-redundancy}

Dropping PC: weaken the sequence policy so it only examines the acting
principal's local chain history, not the global history. Two chains, each
individually compliant, can then produce a forbidden sequence.

\begin{verbatim}
$ python3 redteam/pc/alphabet_escape.py
[PC] ATTACK SUCCEEDED - CI2 violated
[OK] PC: alphabet escape
\end{verbatim}

The TLC trace shows a history ending in two consecutive \texttt{send} actions.
CI2 is violated.

\subsection{RS non-redundancy}

Dropping RS: remove the guard-level revocation check. A revoked principal can
then continue to act in a domain that has not yet received the revocation.

\begin{verbatim}
$ python3 redteam/rs/revocation_race.py
[RS] ATTACK SUCCEEDED - CI3 violated
[OK] RS: revocation race
\end{verbatim}

The TLC trace shows a history entry with $h[i].\mathit{time} \ge
r[h[i].\mathit{principal}]$. CI3 is violated.

This attack also serves as a \emph{mutation test} of CI3: it demonstrates that
CI3 is not vacuously satisfied. Removing the guard that enforces CI3 produces
a violation, so CI3 is genuinely load-bearing.

\subsection{CBAT non-redundancy}

Dropping CBAT: weaken the reachability check to allow any principal to reach
any domain, regardless of attestation. A principal from a chain rooted in
$T_2$ can then act in $T_1$, which $T_2$ does not attest.

\begin{verbatim}
$ python3 redteam/cbat/attestation_laundering.py
[CBAT] ATTACK SUCCEEDED - CI4 violated
[OK] CBAT: attestation laundering
\end{verbatim}

The TLC trace shows a state where the current principal's root domain does
not reach the current domain. CI4 is violated.

\subsection{Summary}

\begin{table}[t]
\centering
\small
\begin{tabular}{@{}llll@{}}
\toprule
\textbf{Attack} & \textbf{Condition dropped} & \textbf{Invariant violated} & \textbf{Result} \\
\midrule
\texttt{pc/alphabet\_escape.py} & PC & CI2 (Sequence Soundness) & succeeds \\
\texttt{rs/revocation\_race.py} & RS & CI3 (Revocation Freshness) & succeeds \\
\texttt{cbat/attestation\_laundering.py} & CBAT & CI4 (Attestation Soundness) & succeeds \\
\bottomrule
\end{tabular}
\caption{Independence of the three conditions. Each attack succeeds, so no
condition can be dropped.}
\label{tab:independence}
\end{table}

This is an \emph{independence} result, not a \emph{necessity} result. It shows
the conditions are non-redundant in this construction. It does not show that
no weaker condition set would suffice, and we do not claim that.

%------------------------------------------------------------ Related work --
\section{Related work}
\label{sec:related}

\subsection{Delegation and access control}

The speaks-for relation and its transitivity are from Lampson, Abadi, Burrows,
and Wobber~\cite{lampson1992authentication}, \S3.2. The same paper introduces
the handoff axiom (P10) in \S3.3 and notes that the general form is ``too
powerful'' for some settings, suggesting a qualified variant---this is the
mechanism our CBAT condition restates. Abadi, Burrows, Lampson, and
Plotkin~\cite{ablp1993calculus} formalize the relation into a calculus for
access control in distributed systems.

SPKI/SDSI~\cite{rfc2693spki} provides a certificate-chain theory for
authorization, with 5-tuple reduction in \S6.3 and validity intervals in
\S5.4. RT (Li, Mitchell, and Winsborough~\cite{rt2002}) introduces linked roles
as a mechanism for bounding transitivity across domains---the closest formal
analog to CBAT's ``explicitly bounded'' clause.

More recently, SentinelAgent~\cite{sentinelagent2026} provides a delegation
chain calculus for AI agents with TLA+ verification. Our delegation layer
restates the standard model; the composition with sequence policies and
federation is the point of departure.

\subsection{Sequence policies and runtime enforcement}

Runtime enforcement of security policies is
Schneider~\cite{schneider2000enforceable}, and the safety/liveness dichotomy
is Alpern and Schneider~\cite{alpern1985defining}. Ligatti, Bauer, and
Walker~\cite{ligatti2005edit} extend enforcement to edit automata, which can
modify actions rather than only block them. Our sequence policies are
safety-only and prefix-closed, so Schneider's framework suffices.

The specific application of sequence policies to agent actions is from
Bounded Agents~\cite{boundedagents2026} (the Agentic Principal Chain). PCAS
provides a Datalog-derived policy compiler for secure agentic
systems~\cite{pacs2026}. Our PC condition is a specific constraint on the
alphabet of the sequence policy; it is not a full policy language.

\subsection{Federation and cross-domain trust}

Federated governance without a trusted center is discussed in Wang's testbed
work~\cite{wang2026federated}, which addresses multi-principal agents that do
not share a central controller. COA-MAS v2~\cite{coamas2026} is a
meta-framework for cross-domain governance, explicitly rejecting a single
federated architecture for all contexts. Our federation graph is a
specific instance of the general problem these works address; we do not
claim to solve the general cross-domain governance problem.

\subsection{Composition of security properties}

The compositionality of safety properties is McLean~\cite{mclean1994general}.
Universally composable security is Canetti~\cite{canetti2001uc}. The
composition of cryptographic protocols is addressed by Abadi and
Gordon~\cite{abadi1997spi}.

Two recent results make the negative side of our problem concrete.
Spera~\cite{spera2026noncompositionality} proves a non-compositionality
theorem: two individually safe agents can jointly reach a forbidden goal
through an emergent conjunctive dependency. AgentRFC~\cite{agentrfc2026}
documents 20 cross-protocol violations across five composition patterns.

SynectivSec is the positive counterpart: given three specific compatibility
conditions, composition preserves the invariants. Where Spera and AgentRFC
establish the negative case, we characterize a positive case.

\subsection{Agent-specific security frameworks}

OAP (Open Agent Passport)~\cite{oap2026} is an open specification for
pre-action authorization of AI tool calls. PCAS~\cite{pacs2026} provides
deterministic policy enforcement for secure agentic systems.
AgentGuardian~\cite{agentguardian2026} learns context-aware access-control
policies from execution traces. These systems operate at the per-action
level; SynectivSec composes above them and does not replace them.

\subsection{Summary}

The individual primitives---delegation chains, sequence policies, federated
attestation---have been studied extensively. What has not been stated is the
\emph{composition} of all three under a preservation theorem. That is the
contribution of this paper.

%------------------------------------------------------------ Limitations --
\section{Limitations and threats to validity}
\label{sec:limitations}

\subsection*{Scope of the proof}

The proof of \Cref{thm:preservation} is by hand. It is not mechanized in a
proof assistant (Lean, Coq, Isabelle, TLAPS). A mechanized proof would
strengthen the claim and is future work.

The proof assumes sequential transitions. A concurrent model with interleaved
transitions would require a refinement argument not given here.

\subsection*{Scope of the verification}

The TLA+ verification is on a bounded instance: $\mathit{time} \le 4$,
$\mathit{Len}(\mathit{history}) \le 3$, 6 principals, 2 chains, 3 domains.
Attempts to scale to larger instances (12 principals, or 6 principals with 4
domains and 4 capabilities) exceeded 5 minutes of runtime due to state-space
explosion. Scaling remains open and would require symmetry reduction, larger
compute, or symbolic model checking (e.g., Apalache).

\subsection*{CI5 non-triviality}

In the current model, $\mathsf{DomainPolicy}$ is a total function, so CI5 is
trivially true. A model with layered or uncertain policies would make CI5
non-trivial. Extending the model to such policies is future work.

\subsection*{Dynamic federation}

The attestation relation $\mathcal{A}$ is a compile-time constant. Domains
joining or leaving the federation during a composition would require a
temporal extension of the theorem.

\subsection*{Independence vs. necessity}

The attack harness establishes independence: each condition is non-redundant
in this construction. It does not establish necessity in the strong sense
(that no weaker set of conditions suffices). We do not claim the stronger
result.

\subsection*{Prior-art verification}

CBAT is verified against Lampson et al.~\cite{lampson1992authentication}
\S3.2 and \S3.3 from the primary source. RS is verified against SPKI
RFC~2693~\cite{rfc2693spki} \S5.4 from the primary source. PC's lineage to
Bounded Agents~\cite{boundedagents2026} and Schneider~\cite{schneider2000enforceable}
is verified from secondary sources; a primary-source check is future work.
The claim that the \emph{joint} formulation of RS is novel has not been
verified by systematic search of the 1990s--2010s trust-management literature.

\subsection*{No reference implementation}

There is no compiled, deployable software. The TLA+ model and the attack
harness are research artifacts. A reference implementation and integration
with agent frameworks (MCP, A2A) is future work.

\subsection*{Non-goals}

SynectivSec addresses action-level and sequence-level constraints. It does not
address content-level attacks (prompt injection, output manipulation),
availability, key compromise, or alignment. These are outside the scope of the
theorem and are handled by other layers of the stack.

%------------------------------------------------------------- Conclusion --
\section{Conclusion}
\label{sec:conclusion}

We have presented a composition theorem for the joint treatment of delegation
chains, sequence policies, and cross-domain federation. The theorem states
three compatibility conditions (PC, RS, CBAT) that are sufficient for the
composition to preserve five safety invariants (CI1--CI5). We prove the
theorem by structural induction over the transition relation, verify it in
TLA+ on a bounded instance with zero violations across 300{,}447 distinct
states, and demonstrate that each condition is non-redundant by a runnable
counterexample.

The individual conditions are not novel. The composition is. The contribution
is the joint treatment, the preservation theorem, and the independence result.

The artifact is released at
\url{https://github.com/kakashi-kx/synectivsec} under Apache~2.0 (code) and
CC-BY~4.0 (documentation). Contributions---especially counterexamples to the
theorem---are welcome.

%---------------------------------------------------------- Bibliography --
\begin{thebibliography}{99}

\bibitem{lampson1992authentication}
B.~Lampson, M.~Abadi, M.~Burrows, and E.~Wobber.
\newblock Authentication in distributed systems: Theory and practice.
\newblock \emph{ACM Transactions on Computer Systems}, 10(4):265--310, 1992.

\bibitem{ablp1993calculus}
M.~Abadi, M.~Burrows, B.~Lampson, and G.~Plotkin.
\newblock A calculus for access control in distributed systems.
\newblock \emph{ACM Transactions on Programming Languages and Systems},
15(4):706--734, 1993.

\bibitem{rfc2693spki}
C.~Ellison, B.~Frantz, B.~Lampson, R.~Rivest, B.~Thomas, and T.~Ylonen.
\newblock SPKI certificate theory.
\newblock RFC 2693, Internet Engineering Task Force, September 1999.

\bibitem{rt2002}
N.~Li, J.~C.~Mitchell, and W.~H.~Winsborough.
\newblock Design of a role-based trust-management framework.
\newblock In \emph{Proc. IEEE Symposium on Security and Privacy}, 2002.

\bibitem{schneider2000enforceable}
F.~B.~Schneider.
\newblock Enforceable security policies.
\newblock \emph{ACM Transactions on Information and System Security},
3(1):30--50, 2000.

\bibitem{alpern1985defining}
B.~Alpern and F.~B.~Schneider.
\newblock Defining liveness.
\newblock \emph{Information Processing Letters}, 21(4):181--185, 1985.

\bibitem{ligatti2005edit}
J.~Ligatti, L.~Bauer, and D.~Walker.
\newblock Edit automata: Enforcement mechanisms for run-time security
policies.
\newblock \emph{International Journal of Information Security},
4(1--2):2--16, 2005.

\bibitem{mclean1994general}
J.~McLean.
\newblock A general theory of composition for trace sets closed under
selective interleaving functions.
\newblock In \emph{Proc. IEEE Symposium on Security and Privacy}, 1994.

\bibitem{canetti2001uc}
R.~Canetti.
\newblock Universally composable security: A new paradigm for cryptographic
protocols.
\newblock In \emph{Proc. IEEE Symposium on Foundations of Computer Science},
2001.

\bibitem{abadi1997spi}
M.~Abadi and A.~D.~Gordon.
\newblock A calculus for cryptographic protocols: The spi calculus.
\newblock In \emph{Proc. ACM Conference on Computer and Communications
Security}, 1997.

\bibitem{sentinelagent2026}
Anonymous.
\newblock SentinelAgent: A delegation chain calculus for AI agents.
\newblock arXiv:2604.02767, 2026.

\bibitem{boundedagents2026}
Anonymous.
\newblock Bounded agents: The agentic principal chain.
\newblock arXiv:2608.15888, 2026.

\bibitem{spera2026noncompositionality}
A.~Spera.
\newblock Non-compositionality of safety for AI agents.
\newblock arXiv:2603.15973, 2026.

\bibitem{agentrfc2026}
Anonymous.
\newblock AgentRFC: Composition safety for agent protocols.
\newblock arXiv:2603.23801, 2026.

\bibitem{oap2026}
U.~Uchibeke.
\newblock Open Agent Passport: Pre-action authorization for AI tool calls.
\newblock arXiv:2603.20953, 2026.

\bibitem{pacs2026}
N.~Palumbo, S.~Choudhary, J.~Choi, P.~Chalasani, and S.~Jha.
\newblock Policy compiler for secure agentic systems.
\newblock arXiv:2602.16708, 2026.

\bibitem{agentguardian2026}
N.~Abaev et al.
\newblock AgentGuardian: Learning access control policies to govern AI agent
behavior.
\newblock arXiv:2601.10440, 2026.

\bibitem{wang2026federated}
K.~Wang.
\newblock A federated governance testbed for multi-principal agents.
\newblock 2026.

\bibitem{coamas2026}
Anonymous.
\newblock COA-MAS v2: A meta-framework for cross-domain agent governance.
\newblock 2026.

\end{thebibliography}

%------------------------------------------------------------ Appendices --
\appendix

\section{Full proof of the preservation theorem}
\label{app:proof}

We give the full structural induction. The notation follows \S\ref{sec:model}.

\subsection{Base case}

Let $s_0 = \langle d_0, \varepsilon, T_0, r_0, 0 \rangle$.

\textbf{CI1.} $|h| = 0$, so the universal quantifier is vacuous.

\textbf{CI2.} $\varphi$ is prefix-closed, and $\varepsilon$ is a prefix of
every history, so $\varepsilon \models \varphi$.

\textbf{CI3.} $|h| = 0$, vacuous.

\textbf{CI4.} $T_0 = \mathsf{RootDomain}(p_0)$ by definition of $s_0$, and
$T_0 \in \mathsf{Reach}(F, T_0)$ by the reflexive base case.

\textbf{CI5.} $\mathsf{DomainPolicy}$ is total by definition.

\subsection{Inductive step: T-Action}

$s \to s'$ where $s' = \langle d, h \frown \langle \langle p, a, t \rangle
\rangle, T, r, t + 1 \rangle$, subject to G1--G5.

\textbf{CI1.} For $i \le |h|$, $h'[i] = h[i]$, so $\mathsf{CI1}(s)$ applies.
For $i = |h| + 1$, $h'[i].\mathit{action} = a \in \kappa(p)$ by G2. So CI1
holds.

\textbf{CI2.} $h' = h \frown \langle a \rangle \models \varphi$ by G3.

\textbf{CI3.} For $i \le |h|$, $h'[i].\mathit{time} = h[i].\mathit{time} <
r[h[i].\mathit{principal}]$ by $\mathsf{CI3}(s)$ and $r' = r$. For
$i = |h| + 1$, $h'[i].\mathit{time} = t < r[p]$ by G2, and
$h'[i].\mathit{principal} = p$, so the inequality holds.

\textbf{CI4.} $T \in \mathsf{Reach}(F, \mathsf{RootDomain}(p))$ by G1, and
$p$ is now the current principal.

\textbf{CI5.} $T$ and the historical actions are unchanged; $a \in \Sigma$ by
PC; $\mathsf{DomainPolicy}$ is total.

\subsection{Inductive step: T-Revoke}

$s \to s'$ where $s' = \langle d, h, T, r[p \mapsto t], t + 1 \rangle$,
subject to $r[p] = \infty$.

\textbf{CI1, CI2, CI4, CI5.} History and current domain unchanged; immediate.

\textbf{CI3.} Let $i \le |h|$. If $h[i].\mathit{principal} \neq p$, then
$r'[h[i].\mathit{principal}] = r[h[i].\mathit{principal}]$ and $\mathsf{CI3}(s)$
applies. If $h[i].\mathit{principal} = p$, we have $h[i].\mathit{time} < t$
because actions are recorded at strictly earlier times than $t$, so
$h[i].\mathit{time} < t = r'[p]$.

\subsection{Inductive step: T-Cross}

$s \to s'$ where $s' = \langle p, h, T', r, t + 1 \rangle$, subject to
$T' \in \mathsf{Reach}(F, \mathsf{RootDomain}(p))$.

\textbf{CI1, CI2, CI3.} History unchanged.

\textbf{CI4.} $T' \in \mathsf{Reach}(F, \mathsf{RootDomain}(p))$ by the guard,
and $p$ is now the current principal.

\textbf{CI5.} Unchanged.

\subsection{Conclusion}

By structural induction, every reachable state $s \in \mathsf{Reach}(D \otimes
S \otimes F)$ satisfies CI1--CI5. \qed

\section{TLA+ module}
\label{app:tla}

The complete TLA+ module is in the artifact at
\texttt{tla/Composition.tla}. We reproduce the invariant definitions and the
transition relation here for reference.

\begin{lstlisting}[language=TLA+]
DoAction(p, a, dom) ==
    /\ ChainRootCanReachDomain(p, dom)
    /\ ActionPermitted(p, a, dom)
    /\ Len(history) < MaxHistory
    /\ SequencePolicyOK(Append(history,
         [principal |-> p, action |-> a, at |-> time]))
    /\ currentPrincipal' = p
    /\ history' = Append(history,
         [principal |-> p, action |-> a, at |-> time])
    /\ currentDomain' = dom
    /\ revokedAt' = revokedAt
    /\ time' = time + 1

Revoke(p) ==
    /\ revokedAt[p] = NeverRevoked
    /\ revokedAt' = [revokedAt EXCEPT ![p] = time]
    /\ currentPrincipal' = currentPrincipal
    /\ history' = history
    /\ currentDomain' = currentDomain
    /\ time' = time + 1

CrossDomain(p, dom) ==
    /\ ChainRootCanReachDomain(p, dom)
    /\ currentPrincipal' = p
    /\ history' = history
    /\ currentDomain' = dom
    /\ revokedAt' = revokedAt
    /\ time' = time + 1
\end{lstlisting}

The full module includes the capability lattice, the attestation relation, the
sequence policy, and the five invariants.

\end{document}
